\section{How Composition Thinks}
\label{sec:ch06-how-composition-thinks}
\index{composition|(}%

Composition is the step between writing subgraph schemas and running a router.
It takes every subgraph's schema, checks them against each other, and either
produces a single supergraph schema or refuses with an error that names the
problem and the location.

The algorithm is chapter~\ref{ch:composition}'s subject, and it runs in a CLI
tool and in the router and in CI, often with different settings. This section is
the model behind the tool: what the composer checks, why it rejects what it
does, and what the errors are actually saying.

\subsection*{The rule that drives everything else}

\index{composition!satisfiability}%
A supergraph schema has to be
satisfiable~\autocite{apollo2026compositionrules}: for every valid query a client could
write, every field in the selection set must be resolvable by at least one
subgraph. That is the one rule from which the others follow, and it is harder
than it looks because the query that breaks it might never have been written by
a human. A client that asks for \code{Product.reviews} gets a response from
Reviews, and a client that asks for \code{Product.weight} gets a response from
Catalog, and a client that asks for both in one query gets both. The question
the composer asks is: does any plausible traversal of this schema lead to a
field nobody can resolve?

If Catalog defines \code{Product.price} and Inventory does not, and a query
starts in Inventory, the router can fetch \code{Product} from Catalog for the
\code{price}. That is satisfiable. If Catalog does not define \code{price}
either, and nothing else does, the field exists in the supergraph and nobody can
resolve it. That is unsatisfiable, and composition refuses.

\subsection*{The explicit rules}

\index{composition!rules}%
The satisfiability rule is the principle, and the compiler's checks are the
machinery. The Apollo Federation specification lists them, and these are the
ones that produce the most errors during development:

\index{composition!duplicate field}%
\medskip\noindent\textbf{One field, one owner, unless shared.} A field of an
object type may appear in exactly one subgraph, or in several if every
occurrence is marked \code{@shareable}. A field that appears twice without
\code{@shareable} on both is a composition error, and the message names the type
and the conflicting subgraphs. This is the error most people hit first.

\index{composition!incompatible types}%
\medskip\noindent\textbf{Shared fields must agree on type.} Two subgraphs
declaring \code{Product.title} both as \code{String!} is fine. One declaring
\code{String!} and the other \code{String} is not. The types must be identical,
nullability and list wrappers included. The composition algorithm compares the
fully resolved type, not the SDL text, so \code{[String]!} and
\code{[String!]!} are different.

\index{composition!external reachability@\code{@external} reachability}%
\medskip\noindent\textbf{\code{@external} fields must be reachable.} A field
marked \code{@external} in one subgraph must be resolvable by another in every
query that could reach it. If the field is in a \code{@requires} selection, the
composer checks that the route from the entity to that field exists.

\index{composition!override validity@\code{@override} validity}%
\medskip\noindent\textbf{\code{@override} must target a valid owner.} The
\code{from} argument must name a subgraph that actually defined the field before
the override. A typo in \code{from} is a composition error, and so is overriding
a field that was never defined.

\index{composition!key matching@\code{@key} matching}%
\medskip\noindent\textbf{Keys must match across subgraphs.} When Catalog defines
\code{Product @key(fields: "id")} and Reviews also defines
\code{Product @key(fields: "id")}, the keys must name the same fields of
compatible types. A key of \code{fields: "sku"} in one and \code{fields: "id"}
in the other means the entity has two separate keys, which is valid; a key of
\code{fields: "id"} of type \code{ID!} in one and \code{String!} in the other is
not.

\index{composition!interfaceObject rules@\code{@interfaceObject} rules}%
\medskip\noindent\textbf{\code{@interfaceObject} must resolve for every
implementing type.} The subgraph that applies \code{@interfaceObject} must have
a reference resolver for every concrete type implementing the interface, or the
composer cannot guarantee that a query reaching the interface will reach a
resolvable entity.

\subsection*{What composition does not check}

\index{composition!limitations}%
A list of what is outside the composer's scope is as useful as the list of what
is inside it, because everything not on it is a runtime surprise.

\begin{itemize}
\item \textbf{Resolver behaviour.} Composition checks the schema. It cannot know
  whether a resolver returns correct data, crashes, or is slow. The supergraph
  is satisfiable on paper and still broken in production, and the only
  detector is testing.

\item \textbf{Data consistency.} A shareable field is a promise two subgraphs
  agree on its value. Composition validates the types, not the values. If
  Catalog has \code{title: "Table"} and Inventory has \code{title: "Desk"},
  composition passes and the client gets whichever one the router fetched.
  Chapter~\ref{ch:strangling-the-monolith} is where this stops being
  theoretical.

\item \textbf{Performance.} A satisfiable query plan can still be slow.
  \code{@requires} that pulls a deep tree, a chain of entities six subgraphs
  long, a reference resolver without a DataLoader: the composer signs off
  on all of them. Chapter~\ref{ch:resilience-and-performance} has the
  profiling.

\item \textbf{Cycles.} Two subgraphs can each \code{@require} a field from the
  other, and composition accepts the arrangement. At runtime the router detects
  the cycle and returns an error, which is a worse position than having the
  composer say no before deployment. Chapter~\ref{ch:hard-modeling-problems}
  covers how to spot them before the schema is committed.
\end{itemize}

\subsection*{Reading a composition error}

\index{composition!error messages}%
The first time composition fails, the message looks like a stack trace from a
language you did not write. Here is a real one, paraphrased:

\begin{minted}{text}
SATISFIABILITY_ERROR: The following supergraph API query:
  { product(sku: "MOS-FRN-0001") { reviews { author { displayName } } } }
cannot be satisfied by the subgraphs because:
  - Customer.displayName is not resolvable from subgraph "reviews"
  - Customer is not defined in subgraph "reviews"
\end{minted}

The error names a query, not just a type and a field. That is the key to reading
it: the composer found a traversal path that reaches a field nobody owns, and it
printed the path. The fix is either to mark the field \code{@shareable} where it
is, to remove it from where it should not be, or to add \code{@external} and a
\code{@requires} that tells the router where to get it.

Chapter~\ref{ch:composition} and appendix~\ref{app:composition-errors} decode
each error class with the fix and the example. For now, the point is that a
composition error is answering a question you want answered before deployment:
does this schema work? The answer is sometimes no, and getting it from a
compiler is better than getting it from a client at runtime.

\index{composition|)}%
